% ============ 第1章 §1.9 Hilbert变换的极大算子与主值的逐点收敛(原书 pp.70-78) ============
\section{Hilbert 变换的极大算子与主值的逐点收敛}\label{sec:1.9}
\index{Hilbert 变换!极大算子}
\index{Cotlar 不等式}
\index{Cotlar--Knapp--Stein 概正交引理}


当 $1 \leq p < \infty$ 时,若 $f \in L^{p}(\mathbb{R})$,则对一切 $\varepsilon > 0$,
\begin{equation*}
H_{\varepsilon}f(x) := \frac{1}{\pi}\int_{|y - x| > \varepsilon}\frac{1}{x - y}f(y)\dd y
\end{equation*}
存在.这是因为对每个 $x$,$\frac{1}{x-y}\chi_{|x-y|>\varepsilon}(y) \in L^{p'}(\mathbb{R})$,由 Hölder 不等式即可保证.于是,定义 Hilbert 变换的极大算子 $H_{*}$ 为
\begin{equation*}
H_{*}f(x) = \sup_{\varepsilon > 0}|H_{\varepsilon}f(x)|.
\end{equation*}
于是得到称为 Cotlar 不等式的下述不等式.

\begin{theorem}[Cotlar 不等式]\label{thm:1.9.1}
设 $1 \leq p < \infty$,$0 < \delta \leq 1$.则对任意 $f \in L^{p}(\mathbb{R})$,存在常数 $C_{\delta}$ 与 $C_{2}$ 使得
\begin{equation*}
H_{*}f(x) \leq C_{\delta}\left\{\left(M\left(|Hf|^{\delta}\right)(x)\right)^{1/\delta} + C_{2}Mf(x)\right\} \quad (x \in \mathbb{R}).
\end{equation*}
\end{theorem}

\begin{proof}
设 $f \in L^{p}(\mathbb{R})$.首先考虑 $\delta = 1$ 的情形.令 $x = 0$ 不失一般性,只需证明对任意正数 $\varepsilon$,
\begin{equation}\tag{1.9.1}
|H_{\varepsilon}f(0)| \leq C_{1}M(Hf)(0) + C_{2}Mf(0).
\label{eq:1.9.1}\end{equation}
令 $Q = (-\varepsilon/2, \varepsilon/2)$,$\tilde{Q} = (-\varepsilon, \varepsilon)$,$f_{1} = f\chi_{\tilde{Q}}$,$f_{2} = f - f_{1} = f\chi_{\tilde{Q}^{c}}$.则
\begin{equation*}
Hf_{2}(0) = H_{\varepsilon}f(0),
\end{equation*}
且对 $x \in Q$,若 $|y| \geq \varepsilon$,则 $|y - x| \geq |y| - |x| > \max(|y|/2, \varepsilon/2)$,故
\begin{align*}
\pi|Hf_{2}(0) - Hf_{2}(x)|&= \left|\displaystyle\int_{|y| \geq \varepsilon}\frac{f(y)}{-y}\dd y - \int_{|y| \geq \varepsilon}\frac{f(y)}{x - y}\dd y\right|  \\
&\leq \displaystyle\int_{|y| \geq \varepsilon}\left|\frac{1}{-y} - \frac{1}{x - y}\right||f(y)|\dd y = \int_{|y| \geq \varepsilon}\frac{|x|}{|y||x - y|}|f(y)|\dd y \\
&\leq 2\displaystyle\int_{|y| \geq \varepsilon}\frac{|x|}{|y|^{2}}|f(y)|\dd y \leq \varepsilon\sum_{j = 0}^{\infty}\int_{2^{j}\varepsilon \leq |y| < 2^{j+1}\varepsilon}\frac{|f(y)|}{|y|^{2}}\dd y \\
&\leq \displaystyle\sum_{j = 0}^{\infty}\frac{1}{2^{j-2}}\cdot\frac{1}{2^{j+2}\varepsilon}\int_{|y| < 2^{j+1}\varepsilon}|f(y)|\dd y \\
&\leq \displaystyle\sum_{j = 0}^{\infty}\frac{1}{2^{j-2}}Mf(0) = 8Mf(0).
\end{align*}
因此
\begin{equation}\tag{1.9.2}
|H_{\varepsilon}f(0)| \leq |Hf(x)| + |Hf_{1}(x)| + \frac{8}{\pi}Mf(0) \quad (x \in Q).
\label{eq:1.9.2}\end{equation}
若 $H_{\varepsilon}f(0) = 0$,则 \eqref{eq:1.9.1} 成立,故设 $H_{\varepsilon}f(0) \neq 0$.任取满足 $0 < \lambda < |H_{\varepsilon}f(0)|$ 的 $\lambda$.则对每个 $x \in Q$,$|Hf(x)|$,$|Hf_{1}(x)|$,$\frac{8}{\pi}Mf(0)$ 中必有一个大于 $\lambda/3$.因此

(i) $\lambda < \dfrac{24}{\pi}Mf(0)$

或

(ii) $Q \subset \{x \in Q;\ |Hf(x)| > \lambda/3\} \cup \{x \in Q;\ |Hf_{1}(x)| > \lambda/3\}$

二者之一成立.这里由 Chebyshev 不等式,
\begin{equation*}
|\{x \in Q;\ |Hf(x)| > \lambda/3\}| \leq \frac{3}{\lambda}\int_{Q}|Hf(x)|\dd x \leq \frac{3|Q|}{\lambda}M(Hf)(0),
\end{equation*}
由 Hilbert 变换的弱 $(1,1)$ 性,
\begin{align*}
|\{x \in Q;\ |Hf_{1}(x)| > \lambda/3\}| & \leq \dfrac{3C}{\lambda}\| f_{1} \|_{L^{1}} = \dfrac{3C}{\lambda}\int_{|x| < \varepsilon}|f(x)|\dd x \\
& \leq \dfrac{6C|Q|}{\lambda}Mf(0),
\end{align*}
故在 (ii) 的情形,
\begin{equation*}
\lambda < 3M(Hf)(0) + 6CMf(0).
\end{equation*}
(i),(ii) 合并,
\begin{equation*}
\lambda < 3M(Hf)(0) + \left(6C + \frac{24}{\pi}\right)Mf(0).
\end{equation*}
因 $\lambda$ 是满足 $0 < \lambda < |H_{\varepsilon}f(0)|$ 的任意数,令 $\lambda \to |H_{\varepsilon}f(0)|$,
\begin{equation*}
|H_{\varepsilon}f(0)| \leq 3M(Hf)(0) + \left(6C + \frac{24}{\pi}\right)Mf(0).
\end{equation*}

其次转向 $0 < \delta < 1$ 的情形.一般地,$(a + b)^{\delta} < a^{\delta} + b^{\delta}$($a, b > 0$),故由 \eqref{eq:1.9.2},
\begin{equation*}
|H_{\varepsilon}f(0)|^{\delta} \leq |Hf(x)|^{\delta} + |Hf_{1}(x)|^{\delta} + \left(\frac{8}{\pi}Mf(0)\right)^{\delta} \quad (x \in Q).
\end{equation*}
对两边在 $Q$ 上取积分平均后再取 $1/\delta$ 次幂,
\begin{align*}
|H_{\varepsilon}f(0)|&\leq 3^{1/\delta}\left\{\left(\dfrac{1}{|Q|}\displaystyle\int_{Q}|Hf(x)|^{\delta}\dd x\right)^{1/\delta} + \left(\frac{1}{|Q|}\int_{Q}|Hf_{1}(x)|^{\delta}\dd x\right)^{1/\delta}\right. \\
&\quad + \left.\left(\dfrac{1}{|Q|}\displaystyle\int_{Q}\left(\frac{8}{\pi}Mf(0)\right)^{\delta}\dd x\right)^{1/\delta}\right\}  \\
&\leq 3^{1/\delta}\left\{\left(M\left(|Hf|^{\delta}\right)(0)\right)^{1/\delta} + \left(\frac{1}{|Q|}\int_{Q}|Hf_{1}(x)|^{\delta}\dd x\right)^{1/\delta} + \frac{8}{\pi}Mf(0)\right\}.
\end{align*}
最后一行中间的一项,由 Hilbert 变换的弱 $(1,1)$ 性与 Kolmogorov 不等式,
\begin{align*}
\left(\dfrac{1}{|Q|}\displaystyle\int_{Q}|Hf_{1}(x)|^{\delta}\dd x\right)^{1/\delta} & \leq \left(\dfrac{C_{\delta}'}{|Q|}\left(C_{1}'\| f_{1} \|_{L^{1}}\right)^{\delta}|Q|^{1-\delta}\right)^{1/\delta} \\
& \leq (C_{\delta}')^{1/\delta}C_{1}'\dfrac{1}{|Q|}\int_{\tilde{Q}}|f(x)|\dd x \\
& \leq 2(C_{\delta}')^{1/\delta}C_{1}'Mf(0).
\end{align*}
因此
\begin{equation*}
|H_{\varepsilon}f(0)| \leq C_{\delta}\left\{\left(M\left(|Hf|^{\delta}\right)(0)\right)^{1/\delta} + C_{2}Mf(0)\right\},
\end{equation*}
证明完毕.
\end{proof}

\begin{corollary}\label{cor:1.9.2}
(i) 当 $1 < p < \infty$ 时,
\begin{equation*}
\| H_{*}f \|_{L^{p}} \leq C_{p}\| f \|_{L^{p}} \quad (f \in L^{p}(\mathbb{R})).
\end{equation*}

(ii) $p = 1$ 时,
\begin{equation*}
|\{x \in \mathbb{R};\ H_{*}f(x) > \lambda\}| \leq C_{1}\frac{\| f \|_{L^{1}}}{\lambda} \quad (f \in L^{1}(\mathbb{R}),\ \lambda > 0).
\end{equation*}
\end{corollary}

\begin{proof}
(i) $1 < p < \infty$ 的情形.设 $f \in L^{p}(\mathbb{R})$.由 Cotlar 不等式的 $\delta = 1$ 情形,
\begin{equation*}
H_{*}f(x) \leq CM(Hf)(x) + CMf(x),
\end{equation*}
故由 Minkowski 不等式,极大函数的 $L^{p}$ 不等式(定理~\ref{thm:1.5.5})以及 Hilbert 变换的 $L^{p}$ 不等式(定理~\ref{thm:1.7.2}),
\begin{align*}
\| H_{*}f \|_{L^{p}} & \leq C\| M(Hf) \|_{L^{p}} + C\| Mf \|_{L^{p}} \\
& \leq CC_{p}\| Hf \|_{L^{p}} + CC_{p}\| f \|_{L^{p}} \\
& \leq CC_{p}C_{p}'\| f \|_{L^{p}} + CC_{p}\| f \|_{L^{p}} = C''\| f \|_{L^{p}}.
\end{align*}

(ii) $p = 1$ 的情形.设 $f \in L^{1}(\mathbb{R})$.固定 $0 < \delta < \eta < 1$,使用 Cotlar 不等式(为简单起见令 $C_{\delta}C_{2} = C_{3}$),
\begin{equation}\tag{1.9.3}\begin{aligned}
\{x \in \mathbb{R};\ H_{*}f(x) > \lambda\}&\subset \left\{x \in \mathbb{R};\ C_{\delta}\left(M\left[(Hf)^{\delta}\right](x)\right)^{1/\delta} > \frac{\lambda}{2}\right\} \cup \left\{x \in \mathbb{R};\ C_{3}Mf(x) > \frac{\lambda}{2}\right\}.
 \end{aligned}\label{eq:1.9.3}\end{equation}
现在,对非负值 $g \in L_{\mathrm{loc}}^{1}(\mathbb{R})$,由 $M(g^{\delta}\chi_{\{g \leq \lambda/2\}})(x) \leq (\lambda/2)^{\delta}$,
\begin{align*}
& \{x \in \mathbb{R};\ \left(M(g^{\delta})(x)\right)^{1/\delta} > \lambda\} \\
& \qquad \subset \left\{x \in \mathbb{R};\ M(g^{\delta}\chi_{\{g > \lambda/2\}})(x) + M(g^{\delta}\chi_{\{g \leq \lambda/2\}})(x) > \lambda^{\delta}\right\} \\
& \qquad \subset \left\{x \in \mathbb{R};\ M(g^{\delta}\chi_{\{g > \lambda/2\}})(x) > (1 - 1/2^{\delta})\lambda^{\delta}\right\} \\
& \qquad = \left\{x \in \mathbb{R};\ \left(M(g^{\delta}\chi_{\{g > \lambda/2\}})(x)\right)^{\eta/\delta} > (1 - 1/2^{\delta})^{\eta/\delta}\lambda^{\eta}\right\}
\end{align*}
故注意 $\eta/\delta > 1$,依次使用 Chebyshev 不等式与极大函数的 $L^{p}$ 不等式,
\begin{align*}
|\{x \in \mathbb{R};\ \left(M(g^{\delta})(x)\right)^{1/\delta} > \lambda\}|&\leq \dfrac{1}{(1 - 1/2^{\delta})^{\eta/\delta}\lambda^{\eta}}\displaystyle\int_{\mathbb{R}}\left(M(g^{\delta}\chi_{\{g > \lambda/2\}})(x)\right)^{\eta/\delta}\dd x  \\
&\leq \dfrac{C^{\eta/\delta}}{(1 - 1/2^{\delta})^{\eta/\delta}\lambda^{\eta}}\displaystyle\int_{\mathbb{R}}\left(g^{\delta}\chi_{\{g > \lambda/2\}}\right)^{\eta/\delta}(x)\dd x \\
&\leq \dfrac{C}{\lambda^{\eta}}\displaystyle\int_{\{g > \lambda/2\}}g^{\eta}(x)\dd x.
\end{align*}
因此
\begin{equation}\tag{1.9.4}
|\{x \in \mathbb{R};\ C_{\delta}\left(M\left[(Hf)^{\delta}\right](x)\right)^{1/\delta} > \lambda/2\}|
\label{eq:1.9.4}\end{equation}
\begin{equation}\tag{1.9.5}
\leq \frac{C}{\lambda^{\eta}}\int_{\{|Hf| > \lambda/(4C_{\delta})\}}|Hf(x)|^{\eta}\dd x
\label{eq:1.9.5}\end{equation}
\begin{equation}\tag{1.9.6}
\leq \frac{C}{\lambda^{\eta}}\left|\{|Hf| > \lambda/(4C_{\delta})\}\right|^{1-\eta}(C_{H}\| f \|_{L^{1}})^{\eta}
\label{eq:1.9.6}\end{equation}
\begin{equation}\tag{1.9.7}
\leq \frac{C}{\lambda^{\eta}}\left(4C_{\delta}C_{H}\| f \|_{L^{1}}/\lambda\right)^{1-\eta}(C_{H}\| f \|_{L^{1}})^{\eta}
\label{eq:1.9.7}\end{equation}
\begin{equation}\tag{1.9.8}
\leq C_{\delta,\eta,H}\frac{\| f \|_{L^{1}}}{\lambda}.
\label{eq:1.9.8}\end{equation}
从 \eqref{eq:1.9.5} 到 \eqref{eq:1.9.6},注意 $|\{x \in \mathbb{R};\ |Hf(x)| > \lambda\}| \leq C_{H}\| f \|_{L^{1}}/\lambda$,并应用了 Kolmogorov 不等式(定理~\ref{thm:1.8.5}).

另外,由极大函数的弱 $(1,1)$ 性,
\begin{equation}\tag{1.9.9}
\left|\left\{x \in \mathbb{R};\ C_{3}Mf(x) > \frac{\lambda}{2}\right\}\right| \leq \frac{2C_{3}}{\lambda}C'\| f \|_{L^{1}}.
\label{eq:1.9.9}\end{equation}
合并 \eqref{eq:1.9.3},\eqref{eq:1.9.8},\eqref{eq:1.9.9},
\begin{equation*}
|\{x \in \mathbb{R};\ H_{*}f(x) > \lambda\}| \leq C\frac{\| f \|_{L^{1}}}{\lambda}.
\end{equation*}
\end{proof}

由这一事实,$H_{\varepsilon}f(x)$ 到 $Hf(x)$ 的 $L^{p}$ 收敛性与逐点收敛性可分别由下面两个推论证明.

\begin{corollary}\label{cor:1.9.3}
下列命题成立.

(i) 当 $1 < p < \infty$ 时,

(a) $\| H_{\varepsilon}f \|_{L^{p}} \leq C_{p}\| f \|_{L^{p}}$($f \in L^{p}(\mathbb{R}),\ \varepsilon > 0$).

(b) $\| H_{\varepsilon}f - Hf \|_{L^{p}} \to 0$($\varepsilon \to 0$).

(ii) 当 $p = 1$ 时,对任意正数 $\lambda$,
\begin{equation*}
\lim_{\varepsilon \to 0}|\{x \in \mathbb{R};\ |H_{\varepsilon}f(x) - Hf(x)| > \lambda\}| = 0,
\end{equation*}
即当 $\varepsilon \to 0$ 时,$H_{\varepsilon}f$ 依测度收敛于 $Hf$.
\end{corollary}

\begin{proof}
(i) 的 (a) 由 $|H_{\varepsilon}f(x)| \leq H_{*}f(x)$ 与系~\ref{cor:1.9.2} 显然.

(i) 的 (b).设 $f \in L^{p}(\mathbb{R})$.其次任取 $\eta > 0$.则由 $\mathcal{S}(\mathbb{R})$ 在 $L^{p}(\mathbb{R})$ 中稠密,存在 $f_{0} \in \mathcal{S}(\mathbb{R})$ 满足
\begin{equation*}
\| f - f_{0} \|_{L^{p}} < \eta.
\end{equation*}
又因 $f_{0} \in \mathcal{S}(\mathbb{R})$,由引理~\ref{lem:1.3.2},存在正数 $\varepsilon_{0}$ 使当 $0 < \varepsilon < \varepsilon_{0}$ 时
\begin{equation*}
\| H_{\varepsilon}f_{0} - Hf_{0} \|_{L^{p}} < \eta.
\end{equation*}
于是当 $0 < \varepsilon < \varepsilon_{0}$ 时,也利用 Hilbert 变换的 $L^{p}$ 有界性,
\begin{align*}
\| H_{\varepsilon}f - Hf \|_{L^{p}} & \leq \| H_{\varepsilon}f - H_{\varepsilon}f_{0} \|_{L^{p}} + \| H_{\varepsilon}f_{0} - Hf_{0} \|_{L^{p}} + \| Hf_{0} - Hf \|_{L^{p}} \\
& \leq C_{p}\| f - f_{0} \|_{L^{p}} + \eta + C_{p}\| f - f_{0} \|_{L^{p}} \\
& \leq (2C_{p} + 1)\eta.
\end{align*}
这表明
\begin{equation*}
\lim_{\varepsilon \to 0}\| H_{\varepsilon}f - Hf \|_{L^{p}} = 0.
\end{equation*}

(ii) 设 $f \in L^{1}(\mathbb{R})$,固定任意正数 $\lambda$.其次任取正数 $\eta$.然后与 (i) 的 (b) 一样,取 $f_{0} \in \mathcal{S}(\mathbb{R})$ 使
\begin{equation*}
\| f - f_{0} \|_{L^{1}} < \eta.
\end{equation*}
则由引理~\ref{lem:1.3.2},
\begin{equation}\tag{1.9.10}
|H_{\varepsilon}f_{0}(x)|,\ |Hf_{0}(x)| \leq C_{0}/(1 + |x|), \quad \lim_{\varepsilon \to 0}H_{\varepsilon}f_{0}(x) = Hf_{0}(x),
\label{eq:1.9.10}\end{equation}
故由前者,当 $1 + |x| > 6C_{0}/\lambda$ 时 $|H_{\varepsilon}f_{0}(x) - Hf_{0}(x)| \leq 2C_{0}/(1 + |x|) \leq \lambda/3$,故令 $A = 6C_{0}/\lambda - 1$,
\begin{align*}
& \{x \in \mathbb{R};\ |H_{\varepsilon}f_{0}(x) - Hf_{0}(x)| > \lambda/3\} \\
& \qquad \subset \{x \in \mathbb{R};\ |x| \leq A,\ |H_{\varepsilon}f_{0}(x) - Hf_{0}(x)| > \lambda/3\}.
\end{align*}
因此由 Chebyshev 不等式,
\begin{equation*}
\left|\left\{x \in \mathbb{R};\ \left|H_{\varepsilon}f_{0}(x) - Hf_{0}(x)\right| > \lambda/3\right\}\right| \leq \frac{3}{\lambda}\int_{-A}^{A}\left|H_{\varepsilon}f_{0}(x) - Hf_{0}(x)\right|\dd x.
\end{equation*}
对 \eqref{eq:1.9.10} 的右端使用 Lebesgue 收敛定理,可知存在正数 $\varepsilon_{0}$ 使当 $0 < \varepsilon < \varepsilon_{0}$ 时
\begin{equation*}
\left|\left\{x \in \mathbb{R};\ \left|H_{\varepsilon}f_{0}(x) - Hf_{0}(x)\right| > \lambda/3\right\}\right| < \eta.
\end{equation*}
于是再利用 $H$ 与 $H_{*}$ 的弱 $(1,1)$ 性(定理~\ref{thm:1.8.2},系~\ref{cor:1.9.2}),
\begin{align*}
& |\{x \in \mathbb{R};\ |H_{\varepsilon}f(x) - Hf(x)| > \lambda\}| \\
& \quad \leq |\{x \in \mathbb{R};\ |H_{\varepsilon}f(x) - H_{\varepsilon}f_{0}(x)| > \lambda/3\}| \\
& \quad + |\{x \in \mathbb{R};\ |H_{\varepsilon}f_{0}(x) - Hf_{0}(x)| > \lambda/3\}| \\
& \quad + |\{x \in \mathbb{R};\ |Hf_{0}(x) - Hf(x)| > \lambda/3\}| \\
& \quad \leq \frac{3C}{\lambda}\| f - f_{0} \|_{L^{1}} + \eta + \frac{3C}{\lambda}\| f - f_{0} \|_{L^{1}} \\
& \quad \leq (6C/\lambda + 1)\eta.
\end{align*}
因此
\begin{equation*}
\lim_{\varepsilon \to 0}|\{x \in \mathbb{R};\ |H_{\varepsilon}f(x) - Hf(x)| > \lambda\}| = 0.
\end{equation*}
\end{proof}

\begin{corollary}\label{cor:1.9.4}
设 $1 \leq p < \infty$.若 $f \in L^{p}(\mathbb{R})$,则当 $\varepsilon \to 0$ 时
\begin{equation*}
H_{\varepsilon}f(x) \to Hf(x) \quad (\text{a.e. } x \in \mathbb{R}).
\end{equation*}
\end{corollary}

\begin{proof}
设 $f \in L^{p}(\mathbb{R})$.为证明起见,只需对实部与虚部分别证明即可,故从一开始就假定为实值函数.首先,由系~\ref{cor:1.9.2} 得
\begin{equation}\tag{1.9.11}
|\{x \in \mathbb{R};\ H_{*}f(x) > \lambda\}| \leq C\left(\frac{\| f \|_{L^{p}}}{\lambda}\right)^{p} \quad (\lambda > 0).
\label{eq:1.9.11}\end{equation}
为证明 $\lim_{\varepsilon \to 0}H_{\varepsilon}f(x) = Hf(x)$,考虑下面的函数:
\begin{equation*}
\Delta f(x) := \limsup_{\varepsilon \to 0}H_{\varepsilon}f(x) - \liminf_{\varepsilon \to 0}H_{\varepsilon}f(x).
\end{equation*}
其次任取 $\eta > 0$.则由 $\mathcal{S}(\mathbb{R})$ 在 $L^{p}(\mathbb{R})$ 中稠密,存在 $g \in \mathcal{S}(\mathbb{R})$ 满足
\begin{equation*}
\| f - g \|_{L^{p}} < \eta.
\end{equation*}
对 $g \in \mathcal{S}(\mathbb{R})$,已知 $\lim_{\varepsilon \to 0}H_{\varepsilon}g(x) = Hg(x)$($x \in \mathbb{R}$),故
\begin{align*}
\limsup_{\varepsilon \to 0}(H_{\varepsilon}f(x) - H_{\varepsilon}g(x)) &= \limsup_{\varepsilon \to 0}H_{\varepsilon}f(x) \\
\liminf_{\varepsilon \to 0}(H_{\varepsilon}f(x) - H_{\varepsilon}g(x)) &= \liminf_{\varepsilon \to 0}H_{\varepsilon}f(x).
\end{align*}
故 $\Delta(f - g)(x) = \Delta f(x)$.因此由 $|\Delta(f - g)(x)| \leq 2H_{*}(f - g)(x)$,再利用 \eqref{eq:1.9.11},
\begin{align*}
|\{x \in \mathbb{R};\ |\Delta f(x)| > 1/N\}| & = |\{x \in \mathbb{R};\ |\Delta(f - g)(x)| > 1/N\}| \\
& \leq |\{x \in \mathbb{R};\ 2H_{*}(f - g)(x) > 1/N\}| \\
& \leq C(2N)^{p}\| f - g \|_{L^{p}}^{p} = C(2N)^{p}\eta^{p}.
\end{align*}
因 $\eta > 0$ 任意,令 $\eta \to 0$,
\begin{equation*}
|\{x \in \mathbb{R};\ |\Delta f(x)| > 1/N\}| = 0.
\end{equation*}
于是因 $\{x \in \mathbb{R};\ |\Delta f(x)| > 0\} = \bigcup_{N=1}^{\infty}\{x \in \mathbb{R};\ |\Delta f(x)| > 1/N\}$,
\begin{equation*}
|\{x \in \mathbb{R};\ |\Delta f(x)| > 0\}| \leq \sum_{N=1}^{\infty}\left|\{x \in \mathbb{R};\ |\Delta f(x)| > 1/N\}\right| = 0.
\end{equation*}
即对几乎所有的 $x \in \mathbb{R}$,$\lim_{\varepsilon \to 0}H_{\varepsilon}f(x)$ 存在.另一方面,由系~\ref{cor:1.9.3},这个极限函数是 $Hf(x)$,故
\begin{equation*}
\lim_{\varepsilon \to 0}H_{\varepsilon}f(x) = Hf(x) \quad (\text{a.e. } x \in \mathbb{R}).
\end{equation*}
\end{proof}
